\documentclass[11pt]{article}
\usepackage{mathblatt}
\newcommand{\flaechenbild}[8]{%
\begin{tikzpicture}[x=1cm,y=1cm]
\draw[thick] (0,0) rectangle (4.5,-4.5);
\draw (3,0) -- (3,-4.5);
\draw (0,-3) -- (4.5,-3);
\node at (1.5,0.3) {$#1$};
\node at (3.75,0.3) {$#2$};
\node at (-0.35,-1.5) {$#3$};
\node at (-0.35,-3.75) {$#4$};
\node at (1.5,-1.5) {$#5$};
\node at (3.75,-1.5) {$#6$};
\node at (1.5,-3.75) {$#7$};
\node at (3.75,-3.75) {$#8$};
\end{tikzpicture}}
\begin{document}
\blattkopf{Terme und binomische Formeln}{Test}
\weit
\einheitenkopf[e1][1 Klammern auflösen]{Einheit 1 von 5 · Klammern auflösen}
\zweigzeile{Hier lernst du, Klammern aufzulösen · neu oder schon bekannt – je nach Buch · keine P10-Aufgabe · baut auf: Terme zusammenfassen}
\setcounter{aufgabe}{9}
\begin{aufgabe}{Ich kann Klammern auflösen (neu in diesem Jahr).}
\anweisung{Löse die Klammer auf.}
\rechenplatz{2}
\begin{gleichungsraster}[1]
\gl{x + (y + 4)} & \gl{3x + (y - 2)} \\
\gl{3\cdot(x + 5)} & \gl{-3\cdot(x - 4)} \\
\end{gleichungsraster}
\end{aufgabe}
\begin{aufgabe}{Ich erkenne, was vor einer Klammer steht.}
Kreuze an, was vor der Klammer steht.
\begin{teile}
\teil $5 + (x - 2)$ \hfill \kreuz{Plus}\kreuz{Minus}\kreuz{Zahl}
\teil $2x - 4\cdot(x - 3)$ \hfill \kreuz{Plus}\kreuz{Minus}\kreuz{Zahl}
\end{teile}
\end{aufgabe}
\begin{aufgabe}{Ich kann eine Sachaufgabe lösen.}
Ein Garten ist $(x + 6)$\,m lang.

\rechteck[seiten={x+6,x,,},punkte=]{4}{2.5}

\flaechenbild{x}{3}{x}{2}{x^2}{}{}{}

\begin{teile}
\teil Stelle einen Term auf.
\teil Wie viel?
\end{teile}
\end{aufgabe}
\begin{aufgabe}{Ich finde den Fehler.}
\begin{teile}
\teil Ben rechnet: \rechnung{9x - (4x - 6) &= 9x - 4x - 6 \\ &= 5x - 6}
\teil $(y+9)^2 = $\\ \kreuz{$y^2 + 81$}\\ \kreuz{$y^2 + 18y + 81$}
\teil $6x + 10$ \hfill \janein
\end{teile}
\begin{teilezwei}
\tz $2x + y$ für $x = 3$ und $y = 4$ \leerfeld & \tz $21^2$ \leerfeld \\
\tz $(x+9)^2$: \feld{a} \feld{b} & \tz $4$ \\
\end{teilezwei}
\end{aufgabe}
\erg{12}{a) $x + y + 4$ \quad b) $3x + y - 2$}
\begin{abhakseite}
\abhakgruppe{Einheit 1 · Klammern auflösen}
\abhakgruppe{\itshape Plus oder Minus}
\abhak{10}{Ich kann Klammern auflösen (neu in diesem Jahr).}
\end{abhakseite}
\end{document}
